\documentclass[10pt,a4paper]{amsart}
\usepackage[margin=19mm]{geometry}
\usepackage{amsmath,amssymb,mathtools}
\usepackage[scr=rsfs,cal=euler]{mathalfa}
\usepackage{xfrac}
\usepackage[unicode,hidelinks,bookmarks=false]{hyperref}
\newcommand{\ie}{i.e.\ }
\newcommand{\cf}{cf.\ }
\newcommand{\resp}{respectively\ }
\newcommand{\Subsection}{\subsection}
\providecommand{\nobreakdash}{\nobreak\hbox{-}}
\setlength{\emergencystretch}{2em}
\multlinegap0pt
\usepackage{bm}
\usepackage{bbm}
\usepackage{dsfont}
\usepackage{xspace}
\usepackage{cancel}
\usepackage{mathtools}
\usepackage{amsthm}
\makeatletter
\@ifundefined{lemma}{\newtheorem{lemma}{Lemma}}{}
\@ifundefined{theorem}{\newtheorem{theorem}{Theorem}}{}
\makeatother
\title[Proximity and Radius in Outerplanar Graphs with Bounded Face]{Proximity and Radius in Outerplanar Graphs with Bounded Faces}
\author{Peter Dankelmann, Sonwabile Mafunda, Sufiyan Mallu}
\date{Source: arXiv:2508.10077v1 (CC BY 4.0)}
\begin{document}
\maketitle

\noindent\emph{Source and licence.} Proximity and Radius in Outerplanar Graphs with Bounded Faces. Version arXiv:2508.10077v1 (CC BY 4.0), distributed under the Creative Commons Attribution 4.0 International licence, as recorded in the arXiv licence metadata for this exact version and shown on the abstract page https://arxiv.org/abs/2508.10077v1. Full rights notice: licenses/arxiv-2508.10077-CC-BY-4.0.txt.\par\smallskip

\noindent\textit{Editorial scope.} Counted unit: the Theorem whose source label is \texttt{theo:rad in Outerpanar} in the article (source0.tex lines 876--881, source label \texttt{theo:rad in Outerpanar}), delivered together with the article's own complete original proof of it (source0.tex lines 883--949, one \texttt{proof} environment of 67 lines). No mathematics is written, repaired or re-derived: statement and proof are the article's own text line for line. Required context retained uncounted: lemm:face-component-has-at-most-n-2/2 vertices (lines 330--333); theo:proximity in 2-connected outerplana graph (lines 392--398). Every definition, display and label of those lines is retained unchanged. Omitted from this package and not redistributed: 1--329, 334--391, 399--875, 882--882, 950--1107. Nothing inside a retained excerpt is omitted or altered: every retained body is delivered exactly as the article's lines are written, so no omission receipt of any kind accompanies this package. Citations: the retained excerpts carry no citation command at all, so no bibliography entry of the article is transcribed and no reference is redistributed. Reference adaptation: none was needed and none was made; every reference inside the delivered unit resolves inside it, so no unresolved reference or citation is produced. Printed slips kept literal: no printed word, symbol, hypothesis or number of the counted statement or of its proof was altered. Layout and class: the article's own class is \texttt{article} with packages including authblk, babel, amsmath, amsfonts, amsthm, amssymb, pgf, tikz, pgfplots, float, graphicx, tikz-cd, comment, cite; the wrapper is the lane's pinned \texttt{amsart} editorial wrapper at 10pt on A4, which restates the theorem-like environments the retained text needs and adds the article's own macro definitions that the retained text needs (none), with extra wrapper packages (bm, bbm, dsfont, xspace, cancel, mathtools). The delivered numbering is the unit's own. Assets: no retained span contains a figure environment, a graphics-inclusion command, a PDF-page inclusion, a table or a TikZ picture; no image file is copied into this package, the package declares no asset dependency and no image file is redistributed with it. Licence: version arXiv:2508.10077v1 of this article is distributed under the Creative Commons Attribution 4.0 International licence, as recorded in the arXiv licence metadata for this exact version and shown on the abstract page https://arxiv.org/abs/2508.10077v1. A full rights notice is delivered at licenses/arxiv-2508.10077-CC-BY-4.0.txt. Characters: every retained excerpt is pure ASCII, so the delivered engine, XeLaTeX with its default Latin Modern text font, reports no missing character.
\begin{lemma}
\label{lemm:face-component-has-at-most-n-2/2 vertices}
Let $G$ be a $2$-connected outerplane graph of order $n$. Then there exist a face, say $F$, such that every component of $G - V(F)$ has at most $\frac{n-2}{2}$ vertices.
\end{lemma}

\begin{theorem}\label{theo:proximity in 2-connected outerplana graph}
If $G$ is a $2$-connected outerplane graph of order $n$, whose internal faces have lengths at most $q$, then 
\begin{equation*} \label{eq:pi(G) in terms of q,n}
\pi(G) \leq
   \dfrac{n+5}{8}+\dfrac{q^2-4q+9}{8(n-1)}.
\end{equation*}
\end{theorem}

\begin{theorem}\label{theo:rad in Outerpanar}
Let $G$ be a $2$-connected outerplane graph of order $n$ whose faces have length at most 
$\frac{n+2}{4}$. 
Then
\[ {\rm rad}(G) \leq \left\lfloor \frac{n}{4} \right\rfloor + 1. \]
\end{theorem}

\begin{proof}              
Let $F$, $k$, $v_i$, $P_i$ and $p_i$ be as defined in the proof of 
Theorem \ref{theo:proximity in 2-connected outerplana graph}.
As in the proof of Theorem \ref{theo:proximity in 2-connected outerplana graph} we may asssume 
that $p_0$ has the largest value 
among the $p_i$, and that $p_j$ has the largest value among the $p_i$ with $i \neq 0$. 
We may assume without loss of generality that $j \leq \frac{k}{2}$.

Since the vertices of $P_i$ together with $v_i$ and $v_{i+1}$ induce a Hamiltonian graph in 
$G$ that has order $p_i+2$,  every vertex of $P_i$ is within distance 
$\lfloor \frac{p_i+2}{2} \rfloor$ from both, $v_i$ and $v_{i+1}$. Define 
$ \ell := \lfloor \frac{n+4}{4} \rfloor - \lfloor \frac{p_0+2}{2} \rfloor+ 1$. 
Note that $\ell \geq 1$ since $p_0 \leq \frac{n-2}{4}$ by 
Lemma \ref{lemm:face-component-has-at-most-n-2/2 vertices}.
Let $u:=v_{\ell}$ if $\ell \leq j$, and $u:=v_j$ if $\ell >j$.
To prove the theorem it suffices to show that for all $x \in V(G)$, 
\begin{equation}
d_G(u, x) \leq \left\lfloor \frac{n+4}{4} \right\rfloor.
\end{equation}  
{\sc Case 1:} $x \in V(P_0)$. \\
Then $d_G(u,v_1) \leq \ell-1$ and $d_G(v_1,x) \leq \lfloor \frac{p_0+2}{2} \rfloor$. Hence
\[ d(u,x) \leq d_G(u,v_1) + d_G(v_1, x) 
\leq \left(\left\lfloor \frac{n+4}{4} \right\rfloor - \left\lfloor \frac{p_0+2}{2} \right\rfloor \right) + \left\lfloor \frac{p_0+2}{2} \right\rfloor
= \left\lfloor \frac{n+4}{4} \right\rfloor, \]
as desired.\\[1mm]
{\sc Case 2:}  $x \in V(P_j)$. \\
If $u=v_j$, then 
\[ d_G(u,x) = d_G(v_j,x) \leq \left\lfloor \frac{p_j+2}{2} \right\rfloor 
   \leq  \left\lfloor \frac{p_0+2}{2} \right\rfloor
   <  \left\lfloor \frac{n+4}{4} \right\rfloor, \]
with the last inequality holding since $p_0 \leq \frac{n-2}{2}$. 

If $u = v_{\ell}$, then, by $j \leq \lfloor \frac{k}{2}\rfloor$ and 
$d_G(v_j,x) \leq  \lfloor \frac{p_j+2}{2} \rfloor$ we have    
\[  d_G(u,x)  =  d_G(v_{\ell},v_j) + d_G(v_j,x) 
      \leq   \left\lfloor \frac{k}{2} \right\rfloor -\ell + \left\lfloor \frac{p_j+2}{2} \right\rfloor. \]
Hence
\begin{align*}
 d_G(u,x) \leq&\;  \left\lfloor \frac{k}{2} \right\rfloor -\left(\left\lfloor \frac{n+4}{4} \right\rfloor - \left\lfloor \frac{p_0+2}{2} \right\rfloor+ 1 \right)+ \left\lfloor \frac{p_j+2}{2} \right \rfloor\\
 \leq& \; \left\lfloor \frac{k+p_0 + p_j+2}{2} \right\rfloor  - \left\lfloor \frac{n+4}{4} \right\rfloor,
\end{align*}
Since $k+p_0+p_j \leq n$, we obtain
\[ d_G(u,x) \leq \left\lfloor \frac{n+2}{2} \right\rfloor - \left\lfloor \frac{n+4}{4} \right\rfloor
            \leq \left\lfloor \frac{n+4}{4} \right\rfloor, \]
as desired. \\[1mm]
{\sc Case 3:} $x \in V(P_i)$ for some $i \in \{1,\ldots,k-1\} - \{0,j\}$. \\
Since $u$ and $v_i$ are on the cycle $F$, we have $d_G(u,v_i) \leq \frac{k}{2}$. 
Moreover, $d_G(v_i,x) \leq \frac{p_i+2}{2}$. Hence 
\begin{equation} \label{eq:radius-of-outerplanar-1} 
d_G(u,x) \leq d_G(u,v_i) + d_G(v_i,x) \leq \frac{k}{2} +\frac{p_i+2}{2}. 
\end{equation}
We now consider two sub-cases. \\[1mm]
{\sc Case 3.1:} $p_i \geq k$. \\
We claim that $p_i \leq \frac{1}{2}(n-k-p_0)$. Indeed, since $k+p_0+p_j+p_i \leq n$,
we have $p_i + p_j \leq n-k-p_0$ and thus, since $p_i \leq p_j$ by the definition of $p_j$, 
the inequality $p_i \leq \frac{1}{2}(n-k-p_0)$ follows.
Substituting this into \eqref{eq:radius-of-outerplanar-1} yields 
\[ d_G(u,x) \leq \frac{n+k-p_0+4}{4}, \]
and since our asumtion $p_i \geq k$ implies that $p_0 \geq k$, the theorem follows in this case. \\[1mm]
{\sc Case 3.2:} $p_i \leq k-1$. \\
Then \eqref{eq:radius-of-outerplanar-1} in conjunction with $k \leq \frac{n+2}{4}$ yields that 
\[ d_G(u,x) \leq 2k \leq \frac{n+4}{4}, \]
and the theorem holds in this case. \\[1mm]
{\sc Case 4:} $x \in V(F)$. \\
Then $d_G(u,x) \leq \frac{k}{2}$. Since $k \leq \frac{n+2}{4}$, we have
$d_G(u,x) \leq \frac{n+2}{8} < \lfloor \frac{n}{4} \rfloor +1$, and the theorem follows.
\end{proof}

\end{document}
